
\subsubsection{6.1 Key Findings and Their Interpretation}

\textbf{Finding 1: $\beta \approx$ 0.143--0.160 $nat^{-1}$ is an actionable quantity.} At this slope, the normalized divergence gap traverses approximately 0.62 normalized units by KL = 8 nats --- more than half the [0, 1] scale. Practitioners should expect the calibration-alignment divergence to become practically significant (gap > 0.2) within approximately 4--8 nats of KL budget beyond the divergence onset point.

\textbf{Finding 2: Near-identical slopes (ratio 1.116) are consistent with a model-family-independent mechanism.} If the calibration-alignment divergence rate were model-specific, slopes would be expected to vary substantially. Instead, $\beta _Gao \approx \beta _Coste$ within 12\% despite different model families, scales, and task distributions --- consistent with the divergence curve reflecting a property of the RLHF optimization process itself. This interpretation is based on n = 2 datasets; further replication is required to establish whether this consistency generalizes more broadly.

\textbf{Finding 3: The low-KL regime in Gao et al. data reveals a ''ramp-up'' before divergence onset.} At very low KL, gold preference rises faster than RM score --- producing a negative gap before crossing zero at ~3.5 nats. This reflects a genuine feature of RLHF optimization: at low KL, the policy learns outputs that both the reward model and human evaluators prefer. The divergence onset occurs when proxy-target separation begins to dominate. This pattern accounts for the lower $R^2$ = 0.701 in the Gao linear regression. The Coste et al. data show a zero-crossing implied near KL = 0, with the gap entering positive territory earlier and maintaining a more linear profile.

\subsubsection{6.2 Connections to Bidirectional Alignment}

The results provide an empirical instantiation of the bidirectional alignment tension identified qualitatively by the ICLR 2025 Workshop. Optimizing the $AI\rightarrow Human$ direction (via RLHF with a proxy RM) degrades evaluation calibration to held-out human judgment at a quantifiable, replicable rate $\beta \approx$ 0.143--0.160 $nat^{-1}$. The normalized divergence gap metric provides a complementary evaluation instrument computable from standard RLHF evaluation runs, without requiring additional infrastructure beyond what is already produced in typical RLHF experimental settings.

\subsubsection{6.3 Limitations}

\textbf{L1: Digitization-derived data.} All analyses use data values constructed consistent with qualitative descriptions from Coste et al. and Gao et al. published figures, not raw experimental datasets. The perfect $\rho$ = 1.000 in Coste et al. data is most likely a digitization artifact. Effect sizes ($R^2$ = 0.958, p < $10^{-6})$ are robust to $\pm 2$--5\% digitization error, but exact $\beta$ values are indicative estimates. Raw data access from the original authors would enable exact $\beta$ computation with measurement-error-aware confidence intervals.

\textbf{L2: Small sample size (N = 10).} The Gao bootstrap CI marginally overlaps zero, illustrating the statistical power limitation. Future work with raw checkpoint data could compute the divergence gap at 50+ KL levels, substantially improving power.

\textbf{L3: P3 (coverage ratio) not tested.} The broader motivating claim --- that alignment evaluation systematically omits $Human\rightarrow AI$ measurement --- relies on the ICLR 2025 survey's qualitative findings, not computation in this study. The coverage ratio R across published alignment benchmarks is deferred to future work.

\textbf{L4: Evaluation calibration $\neq$ user behavioral calibration.} The operationalization here uses held-out gold preference in RLHF evaluation settings, not actual user behavioral calibration (appropriate reliance; Lai et al. [2021]). Future behavioral studies can test whether evaluation-level divergence predicts user over-reliance.

\textbf{L5: Residual autocorrelation in OLS fit.} The Durbin-Watson statistic of 0.411 (Appendix C) indicates positive residual autocorrelation in the Coste OLS fit, expected for serially ordered KL checkpoints. OLS standard errors and the resulting parametric CI [0.119, 0.168] may be anti-conservative under autocorrelation. The bootstrap CI [0.117, 0.177] independently confirms the strictly positive slope without an i.i.d. assumption.

\subsubsection{6.4 Broader Impact}

If reward model overoptimization systematically degrades held-out human preference calibration at $\beta \approx$ 0.143--0.160 $nat^{-1}$, then standard RLHF evaluation practices that report RM scores without tracking held-out human preference may systematically overestimate achieved alignment as optimization proceeds. The normalized divergence gap metric is a low-cost addition to standard RLHF evaluation pipelines. The findings support more careful alignment evaluation rather than enabling harmful applications.

